\documentclass[../../../include/open-logic-section]{subfiles}

\begin{document}

\olfileid{sth}{z}{sep}
\olsection{بېلونه}

له داسې يو اصل نه پېل کوو چې د ساده ادراک ځاے ونيسي:

\begin{axiom}[د بېلونې شېما] د هر فارمول $\phi(x)$ دپاره دا يو
بديهي اصل دے: د هر $A$ دپاره سټ $\Setabs{x \in A}{\phi(x)}$ شتون لري.
\end{axiom}

پام وکړئ چې دا يو يوازينے بديهي اصل نۀ دے. دا د بديهي اصلونو
يوه \emph{شېما} ده. د بېلونې \emph{بې‌نهايته ډېر} بديهي اصلونه
شته؛ د هر فارمول $\phi(x)$ دپاره يو اصل. شېما په همدې مانا
په لاندې ډول هم ليکل کېدلے شي (او معمولاً همداسې ليکل کېږي):

\begin{defish}
د هر هغه فارمول $\phi(x)$ دپاره چې «$S$» پکښې نۀ وي، دا يو
بديهي اصل دے:
\[
	\forall A \exists S \forall x(x \in S \liff (\phi(x) \land x \in A)).
\]
\end{defish}

د \olref[sth][][]{part} په پېل کښې د ياد شوي دود له مخې،
د بېلونې د اصلونو فارمولونه~$\phi$ پارامترونه لرلے شي.\footnote{د دې
خبرې د مانا د تشريح دپاره هغه بحث وګورئ چې له
\olref[sfr][infinite][induction]{natinductionschema} نه سمدستي وروسته راځي.}

د سټونو هغه تجمعي-تکراري تصور چې بيان کړے مو دے، بېلونه په
نېغه توجيه کوي. د دې د ليدلو دپاره $A$ يو سټ وګرځوئ. نو $A$
په يوه پړاو~$S$ کښې جوړېږي (د \stageshier له مخې). ځکه چې $A$
په پړاو~$S$ کښې جوړ شوے دے، د $A$ ټول غړي له پړاو $S$ نه
مخکې جوړ شوي دي (د \stagesacc له مخې). اوس په ځانګړې توګه
هغه ټول سټونه په پام کښې ونيسئ چې د $A$ غړي دي او $\phi$ هم
پوره کوي؛ څرګنده ده چې دا ټول سټونه هم له پړاو~$S$ نه مخکې
جوړ شوي دي. نو دا هم په سټ $\Setabs{x \in A}{\phi(x)}$ کښې په پړاو~$S$
کښې راټولېږي (د \stagesacc له مخې).

د ساده ادراک په خلاف، دا د رسل له تناقضه ځان ساتي. ځکه چې د سټ
$\Setabs{x}{x \notin
x}$ د شتون ادعا په ساده توګه نۀ شو کولے. بلکې، د يو سټ~$A$
د \emph{راکړل کېدو} په صورت کښې، د سټ
$R_A = \Setabs{x \in A}{x \notin x}$ د شتون ادعا کولے شو.
خو له دې نه يوازې دا ثابتېږي چې $R_A \notin R_A$ او
$R_A \notin A$، او په دې دواړو خبرو کښې هېڅ د اندېښنې شے نۀ شته.

خو بېلونه يوه سمدستي او د پام وړ پايله لري:

\begin{thm}\ollabel{thm:NoUniversalSet}
هېڅ \emph{عمومي} سټ نۀ شته، يعنې $\Setabs{x}{x = x}$ شتون نۀ لري.
\end{thm}

\begin{proof}
د تناقض له لارې د ثبوت دپاره فرض کړئ چې $V$ يو عمومي سټ دے.
نو د بېلونې له مخې $R =
\Setabs{x \in V}{x \notin x} = \Setabs{x}{x \notin x}$ شتون لري،
چې د رسل له تناقض سره ټکر کوي.
\end{proof}

د عمومي سټ نۀ شتون---بلکې د سټونو د مراتبو د لړۍ پرانيستے پاتې
کېدل---د تجمعي-تکراري تصور له تر ټولو بنسټيزو مفکورو څخه دي.
نو دا ليدل ګټور دي چې په شهودي توګه دې پايلې ته له بلې لارې هم
رسېدلے شو. يو عمومي سټ بايد د خپل ځان !!a{element} وي.
خو زموږ د تجمعي-تکراري تصور له مخې، هر سټ د مراتبو په لړۍ کښې
(د لومړي ځل دپاره) په هغه لومړني پړاو کښې راڅرګندېږي چې د هغه
له ټولو !!{element}s نه سمدستي وروسته راځي. له دې نه دا لازمېږي
چې \emph{هېڅ} سټ د خپل ځان غړے نۀ دے. ځکه چې د خپل ځان غړے سټ
بايد د لومړي ځل دپاره له هماغه پړاو نه سمدستي وروسته راڅرګند شي
چې پکښې د لومړي ځل دپاره راڅرګند شوے و، او دا نامعقوله ده.
(په \olref[sth][spine][rank]{defnsetrank} کښې به ووينو چې د سټ
د «رتبې» د مفهوم په کارولو سره دا تشريح څنګه لا دقيقه کولے شو.
خو د دې دپاره به لومړے يو څو نور بديهي اصلونه هم ولرو.)

دلته د بېلونې او د غړو له مخې د برابرۍ يو څو نورې پايلې دي.

\begin{prop}\ollabel{prop:emptyexists}
که کوم سټ شتون ولري، نو $\emptyset$ شتون لري.
\end{prop}

\begin{proof}
که $A$ يو سټ وي، نو $\emptyset = \Setabs{x \in A}{x \neq x}$ د بېلونې له مخې شتون لري.
\end{proof}

\begin{prop}
$A \setminus B$ د هر دوو سټونو $A$ او $B$ دپاره شتون لري.
\end{prop}

\begin{proof}
$A \setminus B = \Setabs{x \in A}{x \notin B}$ د بېلونې له مخې شتون لري.
\end{proof}

دا هم څرګندېږي چې (نږدې) خپل‌سري تقاطعګانې شتون لري:

\begin{prop}\ollabel{prop:intersectionsexist}
که $A \neq \emptyset$ وي، نو $\bigcap A = \Setabs{x}{(\forall y \in A)x \in y}$ شتون لري.
\end{prop}

\begin{proof}
$A \neq \emptyset$ وګرځوئ؛ نو يو $c \in A$ شته. بيا
$\bigcap A =
\Setabs{x}{(\forall y \in A)x \in y} = \Setabs{x \in c}{(\forall y \in
A)x \in y}$ د بېلونې له مخې شتون لري.
\end{proof}

خو شرط $A \neq \emptyset$ ته پام وکړئ؛ ځکه چې $\bigcap
\emptyset$ به د تشې رښتينولۍ له مخې عمومي سټ وي، چې له
\olref{thm:NoUniversalSet} سره ټکر کوي.

\end{document}
